\chapter{Dandruff (Seborrheic Dermatitis of the Scalp)}

\section*{Definition}
Dandruff is a common, chronic scalp condition characterized by flaking of the scalp skin. It affects approximately 50 percent of the population worldwide. It is a mild form of seborrheic dermatitis and is not contagious.

\section*{What Happens in Your Body}
Dandruff results from a combination of factors: colonization by the yeast Malassezia globosa which breaks down sebum into irritating fatty acids; individual susceptibility to Malassezia metabolites; and increased sebum production. The inflammatory response accelerates epidermal cell turnover, causing immature keratinocytes to clump together and shed as visible flakes.

\section*{Causes}
\begin{itemize}
\item Malassezia yeast colonization
\item Increased sebum production (seborrhea)
\item Genetic predisposition
\item Hormonal fluctuations
\item Stress and fatigue
\item Cold, dry weather
\item Infrequent shampooing
\item Neurological conditions (Parkinson disease, HIV)
\item Use of harsh hair products
\end{itemize}

\section*{When to See a Doctor}
\begin{itemize}
\item White or yellowish flakes on the scalp, hair, and shoulders
\item Itchy scalp
\item Dry or greasy scales on the scalp
\item Red patches on the scalp
\item Flaking on other sebaceous areas (eyebrows, nose, ears, chest)
\item No significant hair loss or scarring
\item Flakes worsen in winter and improve in summer
\end{itemize}

\section*{Related Symptoms}
\begin{itemize}
\item Seborrheic dermatitis (more severe, inflammatory form)
\item Scalp itching (pruritus)
\item Scalp redness or irritation
\item Dry skin (xerosis)
\item Psoriasis of the scalp (thicker, silvery plaques)
\item Contact dermatitis
\item Eczema (atopic dermatitis)
\end{itemize}

\section*{Self-Care / Home Management}
\begin{itemize}
\item Shampoo regularly (daily or every other day)
\item Use OTC dandruff shampoos (zinc pyrithione, selenium sulfide, ketoconazole)
\item Alternate dandruff shampoos if one becomes less effective
\item Leave shampoo on the scalp for 5 minutes before rinsing
\item Apply coal tar shampoo for stubborn cases
\item Avoid excessive use of styling products
\item Gentle brushing to remove loose flakes
\end{itemize}

\section*{How Your Doctor Will Evaluate You}
\begin{itemize}
\item Clinical diagnosis based on history and scalp examination
\item Rule out other scalp conditions (psoriasis, tinea capitis)
\item Scalp scraping with KOH preparation to exclude fungal infection
\item Skin biopsy in atypical or treatment-resistant cases
\end{itemize}

\section*{Treatment Options}
\begin{itemize}
\item Medicated shampoos: ketoconazole, zinc pyrithione, selenium sulfide, coal tar
\item Topical corticosteroids for inflammation
\item Topical antifungal creams for non-scalp areas
\item Calcineurin inhibitors for refractory seborrheic dermatitis
\item Salicylic acid shampoos to loosen thick scales
\item UV phototherapy for widespread or severe cases
\end{itemize}
\section*{References}
\begin{thebibliography}{99}
\bibitem{ref1} Pierard-Franchimont C, Xhauflaire-Uhoda E, Pierard GE. "Revisiting Dandruff." Int J Cosmet Sci. 2006;28(5):311--318.
\bibitem{ref2} Ranganathan S, Mukhopadhyay T. "Dandruff: The Most Commercially Exploited Skin Disease." Indian J Dermatol. 2010;55(2):130--134.
\bibitem{ref3} Elewski BE. "Clinical Diagnosis of Common Scalp Disorders." J Investig Dermatol Symp Proc. 2005;10(3):190--193.
\bibitem{ref4} Gupta AK, Bluhm R. "Seborrheic Dermatitis." J Eur Acad Dermatol Venereol. 2004;18(1):13--26.
\bibitem{ref5} James WD, Berger TG, Elston DM. "Andrews' Diseases of the Skin: Clinical Dermatology." 12th ed. Elsevier; 2016.
\end{thebibliography}


\clearpage
